\begin{qfmechanism}[Auto-deleveraging on perpetual venues: last-resort loss socialisation, its trilemma, and the 10 October 2025 cascade]\label{mech:crypto-auto-deleveraging}
\qffield{Statement} Chitra describes auto-deleveraging (ADL) as ``a last-resort loss socialization mechanism for perpetual futures venues'', triggered ``when solvency-preserving liquidations fail'', after insurance funds are exhausted, and proves a trilemma: ``no policy can simultaneously satisfy exchange solvency, revenue, and fairness to traders''. On Hyperliquid on 10 October 2025, ADL closed \$2.1 billion of positions in 12 minutes; against transparent benchmark allocations the production algorithm ``overshot the minimum trader profit haircut'', costing profitable traders an estimated \$45.0 to \$51.7 million, about \$653.6 million of positions, and the comparison suggests Binance ``overutilized ADL far more'' (abstract and sec.~1) \citep[arXiv v3, abstract and sec.~1]{cd-chitra-adl-2025}.

\qffield{Conditions} The excess-haircut estimates depend on the author's benchmark allocations and on reconstructed venue data; the author is at Gauntlet, a risk-management firm for crypto protocols.

\qffield{Practice} To analyse a liquidation cascade, write each venue's waterfall (liquidation, insurance fund, ADL or socialised loss) from its rules, and record which positions ADL closed and at what price; compare with \hyperref[def:crypto-trading-risk-layers]{crypto trading risk layers}, which describes one venue's rules.

\qffield{Pitfalls} The paper's numbers are one event; whether Hyperliquid's later open-interest loss is due to ADL is the author's conjecture (sec.~1).

\qffield{See also} \hyperref[mech:perpetual-funding]{perpetual funding}; \hyperref[mod:liquidity-spirals]{liquidity spirals}.

\qffield{Further reading} \nolinkurl{cd-soska-2021}; \nolinkurl{exp-brunnermeier-pedersen-liquidity}.

\end{qfmechanism}
